% Part: proof-theory
% Chapter: normalization
% Section: introduction

\documentclass[../../../include/open-logic-section]{subfiles}

\begin{document}

\olfileid[hi]{pt}{nor}{int}

\olsection{परिचय}

यदि आपने किसी प्रतिबंधक~$!A \lif !B$ को !!{prove} किया है, तो संभव है कि
\Elim{\lif} का उपयोग करके आप उसे~$!B$ के किसी !!{proof} में लगा सकें। इसके लिए
निश्चय ही~$!A$ का कोई !!a{proof}~$\delta_1$ भी चाहिए। बहुत संभव है कि $!A
\lif !B$ का आपका !!{proof}~$\delta_1$,~\Intro{\lif} पर समाप्त हो, जो खुली
मान्यता~$!A$ से~$!B$ के !!a{proof}~$\delta_2$ को~$!A \lif !B$ के प्रमाण में
बदलता है। यदि ऐसा है, तो आपके अंतिम !!{proof} का रूप है
\[
\AxiomC{$\Discharge{!A}{x}$}
\RightLabel{$\delta_2$}
\DeduceC{$!B$}
\DischargeRule{\Intro{\lif}}{x}
\UnaryInfC{$!A \lif !B$}
\AxiomC{}
\RightLabel{$\delta_2$}
\DeduceC{$!A$}
\RightLabel{\Elim{\lif}}
\BinaryInfC{$!B$}
\DisplayProof
\]
आपकी युक्ति इस अर्थ में दक्ष है कि आपने पहले से उपलब्ध वस्तु ($!A \lif !B$ के
!!a{proof}) का पुनः उपयोग किया, पर आपका !!{proof} दक्ष नहीं है। $!A \lif !B$
को प्रस्तुत करके तुरंत उसका विलोपन करना एक चक्कर जैसा लगता है।~$!B$ को
!!{prove} करने का कोई अधिक सीधा मार्ग होना चाहिए।

वास्तव में ऐसा मार्ग सदा होता है। ऊपर जैसे ``चक्कर''—!!a{introduction} नियम
के तुरंत बाद लगा !!a{elimination} नियम—प्राकृतिक निगमन के !!{proof}s से सदा
हटाए जा सकते हैं। इसे \emph{प्रसामान्यीकरण} कहते हैं और परिणाम एक
\emph{प्रसामान्य} !!{proof} (या \emph{प्रसामान्य रूप में} !!a{proof}) होता है।

अनुक्रम कलन के कट-विलोपन प्रमेय की तरह इस परिणाम का प्रमाण कुछ जटिल है और
बहुत-सी स्थितियों पर विचार माँगता है। चिरसम्मत \Log{N1c} के लिए इसे सिद्ध
करना अंतःप्रज्ञात्मक \Log{N1i} की अपेक्षा कठिन है। जैसे कट-विलोपन दिखाता है
कि अनुक्रम कलन में \Cut{} नियम का उपयोग करके आप उससे अधिक कुछ !!{prove} नहीं
कर सकते जितना उसके बिना, वैसे ही प्रसामान्यीकरण दिखाता है कि चक्करों के उपयोग
से आप उनसे बचते हुए किए जा सकने वाले प्रमाणों से अधिक कुछ सिद्ध नहीं कर सकते।
अन्य समानताएँ भी हैं: किसी परिणाम का प्रसामान्य !!{proof}, उसी परिणाम के
चक्कर वाले सबसे छोटे !!{proof} से बहुत लंबा हो सकता है, ठीक जैसे अनुक्रम कलन
में कट वाला !!{proof}, सबसे छोटे कट-मुक्त !!{proof} से बहुत छोटा हो सकता है।
उप-!!{formula} गुण—कि किसी परिणाम को सदा ऐसे !!{proof} से !!{prove} किया जा
सकता है जिसमें केवल उस परिणाम के उप-!!{formula}s प्रयुक्त हों—प्राकृतिक निगमन
के लिए प्रसामान्यीकरण से वैसे ही मिलता है जैसे अनुक्रम कलन के लिए कट-विलोपन
से।

\end{document}
